\begin{proposition}\label{rl-000L}\label{prop:condition-for-monomorphism}
  If $G$ is any connected group scheme over a field acting on a stack $X$ and \(I_X\to X\) is unramified, the map $X^{G}\to X$ is a monomorphism.
\end{proposition}
\begin{proof}
  Suppose we have a map $h:A\to X\in X(A)$ with two lifts $(f,g:A\to X^{G})$ under the natural map $\epsilon:X^{G}\to X$, i.e. that there is an isomorphism $\alpha:f\Rightarrow g$. Since $X^{G}$ parameterizes splittings of the map $\St_{X,G}\to G_X$ (Lemma~\ref{lem:fixed-stack-parameterizes-splittings}) this is the same as having two sections $s_f$ and $s_g$ of $\alpha_A:h^*\St_{X,G}\to G_A$. Since both \(s_f\) and \(s_g\) are sections, composing with \(\alpha_A\) defines the same map. Therefore they differ by some element of \(h^*I_X\) by exactness in~\eqref{eqn:inertia-stack-sequence}; that is, the image of $s_fs^{-1}_g$ lands in the image of the inclusion $h^*I_X\hookrightarrow h^*\St_{X,G}$ where $s_g^{-1}$ is the map $G_A\to h^*\St_{X,G}$ defined by $x\mapsto (s_g(x))^{-1}$. The image of \(s_fs^{-1}_g\) must be connected because \(G\) is connected, and it must land in the identity component of \(h^*I_X\) since this is a homomorphism of group algebraic spaces. However the identity component consists only of the identity itself since \(I_X \to X\) is unramified, hence \(s_fs^{-1}_g\) is trivial.
\end{proof}
